\documentclass[10pt, aspectratio=169]{beamer}
\usetheme{Madrid}
\usecolortheme{default}

% Paquetes necesarios
\usepackage[utf8]{inputenc}
\usepackage[spanish]{babel}
\usepackage{amsmath, amssymb}
\usepackage{graphicx}
\usepackage{booktabs}
\usepackage{hyperref}
\usepackage{listings}
\usepackage{xcolor}

% Configuración de listings para SQL y Java
\lstset{
    basicstyle=\ttfamily\small,
    keywordstyle=\color{blue},
    commentstyle=\color{green!60!black},
    stringstyle=\color{red},
    breaklines=true,
    showstringspaces=false
}

% Título y datos
\title{Diseño de Base de Datos para un Sistema de Información: Proyecto FleetTrac}
\subtitle{Día 1: Fundamentos, Modelado ER y Entorno Geoespacial}
\author{Dr. Guillermo Monroy}
\institute{UAMC}
\date{\today}

\begin{document}

\begin{frame}
    \titlepage
\end{frame}

\begin{frame}{Contenido del Día 1}
    \tableofcontents
\end{frame}

\section{Introducción a los SIG y el proyecto FleetTrack}
\begin{frame}{¿Qué es un Sistema de Información Geoespacial?}
    \begin{itemize}
        \item Combina datos geográficos (ubicaciones, rutas, áreas) con información alfanumérica.
        \item Permite: visualizar en mapas, consultas espaciales (distancia, intersección), análisis de rutas.
        \item \textbf{Ejemplos:} Waze, Google Maps, aplicaciones de logística, monitoreo de flotas.
    \end{itemize}
    \vspace{0.5cm}
    \textbf{Nuestro proyecto: \textit{FleetTrack}}
    \begin{itemize}
        \item Sistema de monitoreo de vehículos y reporte colaborativo de incidentes.
        \item Tecnologías: \textbf{PostgreSQL + PostGIS}, Spring Boot (Java), Leaflet.js.
    \end{itemize}
\end{frame}

\begin{frame}{Ciclo de vida de una base de datos}
    \begin{enumerate}
        \item \textbf{Requisitos:} ¿Qué necesita el sistema? (ej: guardar ubicaciones de vehículos, reportar incidentes, trazar rutas).
        \item \textbf{Diseño conceptual:} Modelo Entidad-Relación (ER) – independiente del SGBD.
        \item \textbf{Diseño lógico:} Traducción a tablas relacionales + normalización.
        \item \textbf{Diseño físico:} Índices, particionamiento, tipos de datos específicos (ej: geometrías PostGIS).
    \end{enumerate}
    \begin{block}{Hoy nos enfocamos en el paso 2: \alert{Modelo ER}}
        Identificamos entidades, atributos y relaciones para FleetTrack.
    \end{block}
\end{frame}

\section{Modelo Entidad-Relación para FleetTrack}
\begin{frame}{Entidades principales}
    \begin{table}
        \centering
        \begin{tabular}{llp{5cm}}
            \toprule
            \textbf{Entidad} & \textbf{Atributos clave} & \textbf{Observación geoespacial} \\
            \midrule
            \texttt{Vehículo} & id, patente, modelo, año & No tiene geometría directa. \\
            \texttt{Conductor} & id, nombre, licencia & -- \\
            \texttt{Ubicación} & id, timestamp, \alert{coordenadas} & Punto geográfico (lat, lon). \\
            \texttt{Incidente} & id, tipo, descripción, \alert{coordenadas}, estado & Punto. \\
            \texttt{Ruta} & id, nombre, \alert{trayecto} & Línea (secuencia de puntos). \\
            \bottomrule
        \end{tabular}
        \caption{Entidades con atributos relevantes}
    \end{table}
\end{frame}

\begin{frame}{Relaciones y cardinalidades}
    \begin{itemize}
        \item \textbf{Conduce:} Conductor $\longrightarrow$ Vehículo (1:N).  
              Un conductor puede manejar varios vehículos (histórico).
        \item \textbf{Registra:} Vehículo $\longrightarrow$ Ubicación (1:N).  
              Cada vehículo reporta muchas ubicaciones en el tiempo.
        \item \textbf{Reporta:} Conductor $\longrightarrow$ Incidente (N:1).  
              Un incidente puede ser reportado por un solo conductor (simplificado).
        \item \textbf{Sigue:} Vehículo $\longrightarrow$ Ruta (N:N).  
              Un vehículo puede recorrer varias rutas; una ruta es recorrida por muchos vehículos.
    \end{itemize}
    \begin{exampleblock}{Diagrama ER resumido}
        \centering
        \texttt{(Conductor) ---conduce--- (Vehículo) ---registra--- (Ubicación)} \\
        \texttt{(Conductor) ---reporta--- (Incidente)} \\
        \texttt{(Vehículo) ---sigue--- (Ruta)}
    \end{exampleblock}
\end{frame}

\begin{frame}[fragile]{Atributos especiales: geometrías}
    \begin{itemize}
        \item En lugar de guardar \texttt{lat} y \texttt{lon} por separado, usamos un solo atributo de tipo \textbf{geometría}.
        \item \textbf{PostGIS} (extensión de PostgreSQL) ofrece:
            \begin{itemize}
                \item \texttt{GEOMETRY(POINT, 4326)} – punto en coordenadas geográficas (lat/lon).
                \item \texttt{GEOMETRY(LINESTRING, 4326)} – secuencia ordenada de puntos (ruta).
                \item \texttt{GEOMETRY(POLYGON, 4326)} – áreas (ej: zona de restricción).
            \end{itemize}
        \item SRID 4326 = sistema de coordenadas WGS84 (GPS).
    \end{itemize}
    \vspace{0.3cm}

    \textbf{Ejemplo de inserción:}

    \begin{lstlisting}[language=SQL]
INSERT INTO ubicacion (vehiculo_id, coordenadas, timestamp)
VALUES (1, ST_SetSRID(ST_MakePoint(-58.3816, -34.6037), 4326), NOW());
    \end{lstlisting}
\end{frame}

\section{Tipos de datos espaciales básicos}
\begin{frame}{POINT, LINESTRING y POLYGON}
    \begin{columns}
        \column{0.5\textwidth}
        \textbf{Punto}
        \begin{itemize}
            \item Una coordenada (X,Y).
            \item Útil para: ubicaciones de vehículos, incidentes.
        \end{itemize}
        \textbf{Línea}
        \begin{itemize}
            \item Conjunto ordenado de puntos.
            \item Útil para: rutas, trayectorias.
        \end{itemize}
        \column{0.5\textwidth}
        \textbf{Polígono}
        \begin{itemize}
            \item Área cerrada (ej: zona de exclusión).
            \item Opcional en nuestro proyecto.
        \end{itemize}
    \end{columns}
    \vspace{0.5cm}
    \begin{block}{Funciones espaciales útiles (veremos más adelante)}
        \texttt{ST\_Distance}, \texttt{ST\_Within}, \texttt{ST\_Length}, \texttt{ST\_AsGeoJSON}
    \end{block}
\end{frame}

\section{Preparación del entorno práctico}
\begin{frame}[fragile]{Configuración inicial (vista previa)}
    \begin{enumerate}
        \item \textbf{PostgreSQL + PostGIS} con Docker:
        \begin{lstlisting}[basicstyle=\ttfamily\small]
docker run --name postgis -e POSTGRES_PASSWORD=admin123 \
           -p 5432:5432 -d postgis/postgis
        \end{lstlisting}
        \item \textbf{Proyecto Spring Boot} (Spring Initializr):
        \begin{itemize}
            \item Dependencias: Web, Data JPA, PostgreSQL Driver.
        \end{itemize}
        \item \textbf{Entidades JPA} mapeando las tablas (Vehículo, Ubicación, etc.).
        \item \textbf{Configuración:} application.properties con URL, usuario, contraseña.
        \item \textbf{Verificación:} Insertar datos de prueba y consultar con DBeaver.
    \end{enumerate}
    \begin{alertblock}{¡Atención!}
        Para usar PostGIS desde JPA necesitamos agregar el dialecto espacial:
        \texttt{hibernate.dialect=org.hibernate.spatial.dialect.postgis.PostgisDialect}
    \end{alertblock}
\end{frame}

\section{Resumen y tareas}
\begin{frame}{Resumen del día}
    \begin{itemize}
        \item Entendimos qué es un SIG y cómo se aplica a logística.
        \item Aprendimos el ciclo de vida de una BD y nos centramos en el modelo ER.
        \item Identificamos las entidades, relaciones y atributos de \textit{FleetTrack}.
        \item Introdujimos tipos de datos espaciales (Point, LineString) y la extensión PostGIS.
        \item Configuramos el entorno base (Docker + Spring Boot).
    \end{itemize}
    \vspace{0.5cm}
    \textbf{Tarea para casa:}
    \begin{enumerate}
        \item Completar el diagrama ER incluyendo todos los atributos (ej: \texttt{incidente.tipo} podría ser ENUM).
        \item Asegurarse de que el contenedor PostgreSQL+PostGIS funciona correctamente.
        \item Opcional: crear las entidades JPA básicas (Vehículo y Conductor).
    \end{enumerate}
\end{frame}

\begin{frame}{Referencias y recursos}
    \begin{itemize}
        \item Documentación PostGIS: \url{https://postgis.net/documentation/}
        \item Spring Data JPA: \url{https://spring.io/projects/spring-data-jpa}
        \item Hibernate Spatial: \url{https://hibernate.org/orm/spatial/}
        \item Leaflet.js: \url{https://leafletjs.com/}
    \end{itemize}
    \vspace{1cm}
    \centering
    \textbf{¡Manos a la obra!}
\end{frame}

\end{document}